\documentclass[11pt,a4paper]{article}
\usepackage[T1]{fontenc}
\usepackage[utf8]{inputenc}
\usepackage[french]{babel}
\usepackage{lmodern}
\usepackage{amsmath,amssymb}
\usepackage[margin=2.35cm]{geometry}
\usepackage{booktabs,longtable,tabularx}
\usepackage{enumitem}
\usepackage{graphicx}
\usepackage{xcolor}
\usepackage{microtype}
\usepackage{float}
\usepackage{hyperref}

\hypersetup{colorlinks=true,linkcolor=blue!45!black,urlcolor=blue!45!black}
\setlist{nosep}
\newcommand{\code}[1]{\texttt{#1}}

\title{Prompt M4.3Live\\
\large Moteur M4.3 pilotable en direct, à la recherche d'un effet rebond}
\author{Lignée M4.3Live --- succède à M4.3 (\code{m4\_3\_credit\_soc}, lu
comme référence structurelle, pas une dépendance),\\
destiné à une instance Claude Opus 5}
\date{17 août 2026 --- révision 2}

\begin{document}
\maketitle

\begin{quote}
\sloppy
\emph{Destinataire : une instance Claude Opus 5 (Claude Code), à qui ce
document est confié comme spécification de tâche complète et autonome. Ce
n'est pas un brouillon à discuter avant de commencer : les décisions de
fond sont prises ci-dessous ; les points explicitement laissés ouverts
(\S11) sont à trancher et documenter toi-même, pas à retourner à
l'utilisateur avant d'avoir essayé.}

\emph{Version : 17 août 2026, révision 2. Lignée : succède à M4.3
(\code{m4\_3\_credit\_soc/}, lu comme référence structurelle ---
\textbf{pas une dépendance}, \S3.5) mais n'est plus une extension au sens
strict : l'institution de principal elle-même change de nature (\S3).
Convention de dossier \code{m4\_3live\_credit\_soc/} inchangée.}

\emph{\textbf{Provenance de cette révision} : ce document remplace la
révision 1 du 17 août (matin). Il intègre (a) les annotations
manuscrites de l'utilisateur portées sur le PDF de la révision 1 (15
annotations, dont une vide, entre 12h44 et 16h21) et (b) le rapport
\code{rapport\_architecture\_offline\_online.pdf} (17 août,
après-midi, copié dans ce dossier pour pérennité) que ces annotations ont
motivé. \textbf{Une annotation est explicitement dépassée par deux
autres, postérieures de 25 à 70 minutes} : à 12h57 l'utilisateur note
« on reste en arithmétique tout du long » à propos de \code{target\_rule}
figé au lancement --- lecture encore valide sur ce point précis (\S2) ---
mais à 13h33 et 14h09 il redéfinit complètement ce que « arithmétique »
signifie (\S3). Ne pas lire la première annotation comme figeant
l'ancienne formule littérale $(K_\ell-K_b)/2$ en toute circonstance.}
\end{quote}

\tableofcontents

\section*{0. Lectures préalables obligatoires, et pièges à ne pas suivre}
\addcontentsline{toc}{section}{0. Lectures préalables obligatoires, et pièges à ne pas suivre}

Avant d'écrire une ligne de code, lire dans cet ordre :

\begin{itemize}
\item \code{m4\_3\_credit\_soc/m4\_3/model.py} (référence structurelle,
  $\sim$800 lignes, \textbf{pas une dépendance à importer}, \S3.5) ---
  c'est le fichier dont ce prompt cite les numéros de ligne ci-dessous ;
  si le fichier a changé depuis, revérifier les citations, ne pas les
  croire sur parole.
\item \code{m4\_3\_credit\_soc/m4\_3/io.py} et
  \code{m4\_3\_credit\_soc/report/rapport\_final.md} (contexte
  scientifique de M4.3 : pourquoi \code{loan\_events}, pourquoi la règle
  de principal historique était arithmétique).
\item \code{simulation\_lab/contracts.py}, \code{simulation\_lab/jobs.py},
  \code{simulation\_lab/web/app.py}, \code{docs/README\_simulation\_lab.md}
  --- l'outil d'orchestration existant, dont ce prompt étend l'usage sans
  le casser.
\item \code{CLAUDE.md} (racine du dépôt) --- règles de travail du dépôt.
\item \textbf{\code{m4\_3live\_credit\_soc/prompts/
  rapport\_architecture\_offline\_online.pdf}} --- étude commandée par
  l'utilisateur sur le calcul efficace, à l'échelle de millions de
  résolutions, du transfert optimal introduit en \S3. Lecture obligatoire
  avant d'écrire une ligne du \S3 : ce prompt en reprend les équations et
  l'architecture recommandée sans les redémontrer.
\item Optionnel mais utile, pas bloquant : les rapports des autres
  lignées du dépôt (\code{m4\_credit\_soc/}, \code{m4b\_credit\_soc\_mini/},
  \code{m4\_2\_credit\_soc/}, \code{m4\_2b\_credit\_soc/}, chacun avec son
  \code{report/}) donnent une vision du fonctionnement intrinsèque de
  cette classe de modèle au-delà de M4.3 (ex. tension entre taille
  moyenne et taille autarcique du système, effet d'une multiplication
  scalaire des paramètres) --- ne pas redécouvrir ce qui y est déjà
  documenté.
\end{itemize}

\textbf{Deux passages de \code{CLAUDE.md} et de
\code{docs/README\_simulation\_lab.md} sont obsolètes et vont t'induire en
erreur si tu les suis sans vérifier :}

\begin{enumerate}
\item \code{CLAUDE.md} désigne \code{m4b\_credit\_soc\_mini/} comme
  « moteur actif » et donne \code{random.Random(seed)} comme convention
  RNG. \textbf{Faux pour ce travail} : le moteur de référence structurelle
  est \code{m4\_3\_credit\_soc/m4\_3/model.py}, et son RNG est
  \code{numpy.random.default\_rng(seed)} (\code{Simulation.\_\_init\_\_},
  \code{m4\_3/model.py:525}) --- convention à reproduire dans le fork
  indépendant de M4.3Live (\S3.5).
\item \code{docs/README\_simulation\_lab.md} liste
  \code{web/templates/index.html} et donne « stdlib + matplotlib » comme
  dépendances. \textbf{Faux} : les templates actuels sont
  \code{launch.html}/\code{results.html} (voir
  \code{simulation\_lab/web/app.py:50-53}), et numpy/scipy sont
  explicitement autorisés en plus de la stdlib. Pas de nouvelle
  dépendance externe (pas de Flask/FastAPI/websockets tiers) sans accord
  explicite de l'utilisateur --- rester sur
  \code{http.server}/\code{ThreadingHTTPServer}.
\end{enumerate}

Règle de confiance à 95~\% du dépôt (\code{CLAUDE.md}) : si tu n'es pas
certain à 95~\% d'un fait sur ce projet, le signaler explicitement (fait
observé / inférence / hypothèse / incertitude). Citer systématiquement
\code{fichier.py:ligne}. Ne jamais inventer un comportement du moteur non
vérifié dans le code.

\section*{1. Objectif scientifique : l'effet rebond, à observer pour l'instant, pas à expliquer}
\addcontentsline{toc}{section}{1. Objectif scientifique : l'effet rebond, à observer pour l'instant, pas à expliquer}

Le moteur M4.3 produit, par pas de temps, une production agrégée
\code{prod\_tot} (\code{Simulation.series}, calculée \code{model.py:563-568}
: $A \cdot K^{\gamma}$ par entité, sommée --- inchangé quelle que soit
l'institution de principal, \S3). La « puissance d'extraction » d'une
entité est portée par deux paramètres de sa fonction de
production/rendement, $A$ et $\gamma$ --- ce ne sont \textbf{pas} des
leviers interchangeables :

\begin{itemize}
\item \textbf{$A$ est le levier primaire.} Il multiplie linéairement la
  production et le rendement marginal --- une hausse de $A$ augmente la
  production pour toute capitalisation, sans ambiguïté de sens.
\item \textbf{$\gamma$ est un levier secondaire, à interprétation non
  monotone.} \code{Config.\_\_post\_init\_\_} (\code{model.py:86-87})
  contraint $\gamma \in\, ]0,1[$. Pour $K<1$, augmenter $\gamma$
  \emph{diminue} $K^{\gamma}$. Toute conclusion sur un effet rebond via
  $\gamma$ doit signaler cette non-monotonie plutôt que traiter $\gamma$
  comme un simple cadran de puissance.
\end{itemize}

\textbf{Définition de travail de l'effet rebond pour ce programme} : une
augmentation de $A$ (et/ou $\gamma$, avec la réserve ci-dessus) appliquée
à une partie de la population --- via la portée \code{fraction}
\emph{ou} la portée \code{new} du \S2, pas nécessairement un
sous-ensemble figé une fois pour toutes --- se traduit, sur
\code{prod\_tot} agrégé du système entier, par une réponse non
proportionnelle à l'augmentation locale --- voire par une baisse --- par
comparaison à l'état qu'aurait eu le système en poursuivant \textbf{sans}
ce changement dynamique de paramètres (une simulation de contrôle de
même graine, \S7), et par comparaison à une intervention de même
amplitude appliquée à la population entière (portée \code{all}). Ce
n'est pas un critère binaire prêt à l'emploi : à toi de figer une
statistique précise (élasticité de \code{prod\_tot} à l'intervention,
fenêtre d'observation post-intervention, agrégation multi-graines) avant
de lancer la campagne de mesure, et de la documenter.

\textbf{L'objectif de cette itération du programme est d'OBSERVER l'effet
rebond, pas de l'expliquer ni de le justifier} (décision explicite de
l'utilisateur). Trois canaux plausibles de propagation d'une hausse
locale --- pour information, comme grille de lecture, \textbf{pas comme
livrable obligatoire} :

\begin{enumerate}
\item \textbf{Production directe} : les entités boostées produisent
  plus, mécaniquement.
\item \textbf{Accumulation de capital} : la production s'ajoute au
  capital de l'entité (\code{model.py:567},
  \code{population.K[entity] += produced}) avant la phase de marché du
  même pas --- une entité boostée entre dans le marché du crédit avec un
  $K$ plus élevé, ce qui change sa position de prêteuse ou d'emprunteuse.
\item \textbf{Effet sur les échanges} : une entité boostée change le
  calcul du transfert (\S3) et, si la piste du \S3.4 aboutit, du taux,
  dans chaque paire où elle apparaît --- donc indirectement le risque de
  faillite des autres entités.
\end{enumerate}

Si une décomposition sur ces trois canaux ressort naturellement de
l'analyse, la documenter dans le rapport de résultats (\S9) comme
analyse complémentaire est bienvenu --- mais ne pas la traiter comme un
critère de succès, et ne pas retarder la campagne de mesure du \S7 pour
l'obtenir. Le résultat attendu est un verdict d'observation (\S7) : effet
rebond présent / absent / non tranchable avec le budget disponible.

\section*{2. Portées d'intervention : sémantique exacte, trois portées, pas deux}
\addcontentsline{toc}{section}{2. Portées d'intervention : sémantique exacte, trois portées, pas deux}

La demande de l'utilisateur (« changer les variables pour toutes les
entités ou seulement pour les nouvelles ») définit deux portées ; ce
programme en exige une troisième pour que l'expérience du \S7 soit
interprétable. Les trois portées, comme fonctionnalité générale du
système, applicables (en principe) à n'importe quel paramètre numérique
de \code{Config} :

\begin{itemize}
\item \textbf{\code{all} (toutes les entités)} : la nouvelle valeur
  s'applique immédiatement à toutes les entités actuellement vivantes
  \emph{et} à toutes les entités futures --- rétroactif et prospectif.
\item \textbf{\code{new} (nouvelles entités seulement)} : la nouvelle
  valeur ne s'applique qu'aux entités nées après l'intervention ; les
  entités déjà vivantes gardent leur valeur d'origine \textbf{pour
  toujours} --- sémantique de vintage/cohorte technologique (le capital
  « embarque » la technologie de sa date de naissance, terminologie
  standard en économie de la croissance ; le mentionner dans le rapport
  de conception, \S9).
\item \textbf{\code{fraction} (fraction $\varphi$ de la population
  vivante, tirage aléatoire)} : au moment $t_0$ de l'intervention, tirer
  une fraction $\varphi \in\, ]0,1]$ (ou une liste explicite
  d'identifiants, pour le débogage) des entités \textbf{actuellement
  vivantes} et leur assigner la nouvelle valeur, une seule fois.
  \textbf{Ce que deviennent les naissances postérieures à $t_0$ (question
  explicitement posée par l'utilisateur en annotation) : une
  intervention \code{fraction} ne modifie PAS la valeur par défaut
  utilisée aux naissances futures.} Seules les entités tirées à $t_0$
  sont affectées ; les entités nées après $t_0$ reçoivent le défaut
  alors en vigueur (celui d'avant l'intervention, sauf si une
  intervention \code{all}/\code{new} distincte est émise séparément sur
  le même paramètre). Raison de ce choix, à ne pas perdre en
  implémentant : c'est ce qui garde l'intensité de traitement
  \textbf{fixe} au moment de la mesure --- la raison d'être même de la
  portée \code{fraction}. Si \code{fraction} modifiait aussi le défaut,
  elle se comporterait comme \code{new} plus un tirage ponctuel, et
  l'intensité dériverait à nouveau dans le temps.

  \textbf{Portée requise pour l'expérience du \S7} : \code{new} a une
  intensité de traitement qui dérive dans le temps (avec $\lambda$
  constant, la part boostée de la population croît de 0 vers un régime
  permanent au fil des naissances/morts, donc mesurer une élasticité
  contre une intensité de traitement mouvante n'a pas de sens) ;
  \code{fraction} fixe l'intensité au moment de la mesure --- c'est la
  portée utilisée pour le protocole principal, \code{all} et \code{new}
  restent des modes d'exploration secondaires (\S7).
\end{itemize}

\textbf{Où la portée a un sens, et où elle n'en a pas --- à ne pas
deviner en silence :}

\begin{itemize}
\item $A$ et $\gamma$ (par entité, \S3) : les trois portées ont un sens
  et doivent toutes être implémentées. \textbf{Attention} (mise en garde
  explicite de l'utilisateur) : modifier $A$ ou $\gamma$ pour une
  cohorte change son rendement marginal, donc son échelle de capital
  naturelle (l'échelle autarcique déjà documentée dans la lignée
  M4.2/M4.2B --- voir \code{m4\_2\_credit\_soc/} ou
  \code{m4\_2b\_credit\_soc/} si la formule doit être retrouvée). Faire
  varier $A$/$\gamma$ pour une cohorte \textbf{sans} ajuster \code{K0}
  en proportion peut changer la tension du système (taille relative à
  l'échelle autarcique) en même temps que l'extraction, ce qui
  confondrait l'expérience du \S7. Ce prompt ne tranche pas si \code{K0}
  doit être co-varié automatiquement avec une intervention
  $A$/$\gamma$ --- c'est une décision de protocole à prendre et à
  documenter explicitement avant la campagne du \S7 (liste des
  décisions ouvertes, \S11), pas un détail à laisser implicite.
\item \code{K0} (capital de naissance) : n'a de sens qu'au moment de la
  naissance --- il n'existe rien à modifier « rétroactivement » sur une
  entité déjà née. \code{all} et \code{new} sont donc \textbf{identiques
  par construction} pour \code{K0} ; \code{fraction} n'a pas de sens pour
  \code{K0} seul (pas d'implémentation requise, sauf si utilisé
  conjointement avec $A$/$\gamma$ comme au paragraphe précédent).
  Documenter cette dégénérescence dans l'IHM (ne pas la cacher en
  désactivant silencieusement un sélecteur sans explication) et dans le
  rapport de conception.
\item \code{lam}, \code{delta}, \code{sigma}, \code{rho},
  \code{eta\_beta}, \code{eta\_n\_ref} : ce sont des paramètres de
  population/marché, pas des attributs d'entité --- \code{new} doit être
  un alias explicite de \code{all} (documenté comme tel dans l'IHM, pas
  masqué), \code{fraction} n'a pas de sens (pas d'implémentation
  requise). Ceci ne les exclut pas comme axes secondaires de
  l'\textbf{expérience} du \S7 elle-même --- des jeux différents de ces
  paramètres, donc des régimes de population/marché différents, peuvent
  être intéressants à observer pour la question du rebond (remarque de
  l'utilisateur) ; seule l'extension de la sémantique de portée
  fraction/new à ces paramètres est hors mandat.
\item \code{target\_rule} : dans M4.3Live, il n'existe plus qu'une seule
  institution de principal (\S3) --- la bascule arithmétique/géométrique
  de M4.2B/M4.3 n'est pas reprise (M4.3Live est indépendant de ce
  moteur, \S3.5). Le champ peut être conservé dans \code{Config} pour
  usages futurs mais n'a, pour ce prompt, aucune valeur alternative à
  exposer ; \textbf{exclu de l'intervention en direct} pour la même
  raison que pour les autres choix institutionnels figés : rien dans la
  question du rebond (\S1, \S7) n'exige de changer d'institution en
  cours de run.
\end{itemize}

\textbf{Minage/synchronisation avec la boucle de pas} : une intervention
arrive de manière asynchrone (requête HTTP) pendant qu'un thread exécute
\code{Simulation.step()} en boucle. Elle doit être \textbf{mise en file
d'attente} et appliquée en \textbf{tout début du prochain appel à
\code{step()}, avant les naissances} --- jamais pendant un \code{step()}
en cours. L'intervention effectivement appliquée doit être journalisée
avec le $t$ réel auquel elle a pris effet. Le tirage \code{fraction}
doit utiliser le générateur RNG de la simulation (\code{self.rng}) pour
rester déterministe et rejouable --- pas un générateur séparé.
\textbf{La mise en pause du direct ne doit consommer aucun tirage RNG} :
mettre en pause puis reprendre et exécuter $N$ pas doit produire une
trajectoire strictement identique à $N$ pas exécutés sans pause ---
invariant à tester (\S8).

\section*{3. L'institution de principal : maximiser la puissance d'extraction jointe}
\addcontentsline{toc}{section}{3. L'institution de principal : maximiser la puissance d'extraction jointe}

\textbf{Ce point remplace intégralement le \S3 de la révision précédente
de ce prompt.} L'utilisateur abandonne la règle arithmétique littérale
de M4.3 ($q_A=(K_\ell-K_b)/2$, \code{m4\_3\_credit\_soc/m4\_3/model.py:
269-291}) comme institution première de M4.3Live : cette formule n'était
qu'un cas particulier d'un problème plus général, celui du montant à
échanger entre deux entités pour \textbf{maximiser leur puissance de
production jointe immédiatement après l'échange}. Le rapport
\code{rapport\_architecture\_offline\_online.pdf} (\S0, lecture
obligatoire) pose et résout ce problème :

$$\max_{-x \le \delta \le y} a(x+\delta)^{\alpha} + b(y-\delta)^{\beta},
\qquad 0<\alpha,\beta<1,\ a,b>0 \qquad \text{(rapport \S1, éq.1)}$$

avec $x=K_b$ (capital de l'emprunteuse), $y=K_\ell$ (capital de la
prêteuse), $(a,\alpha)=(A_b,\gamma_b)$, $(b,\beta)=(A_\ell,\gamma_\ell)$
(rapport \S2, éq.8). En posant $C=x+y$ (la somme conservée par
l'échange), la solution s'écrit $\delta^{*}=C\lambda^{*}-x$ pour un
certain $\lambda^{*}\in[0,1]$ (rapport \S1, éq.2). C'est le calcul de
$\lambda^{*}$ (ou du capital optimal $h(C)=C\lambda^{*}$ attribué à
l'emprunteuse) qui pose un problème de coût selon le régime.

\subsection*{3.1 Trois régimes, un seul principe institutionnel}

\begin{center}
\begin{tabular}{@{}p{2.1cm} p{3.6cm} p{5.9cm} p{2.9cm}@{}}
\toprule
Régime & Condition & Calcul & Coût \\
\midrule
(a) homogène & $A_\ell=A_b$ \textbf{et} $\gamma_\ell=\gamma_b$ &
  $\delta^{*}=(K_\ell-K_b)/2$, la formule historique de M4.3 &
  trivial \\[0.4em]
(b) exposants égaux & $\gamma_\ell=\gamma_b$, $A_\ell\neq A_b$ possible &
  forme fermée exacte, éq.~(\ref{eq:lambda-star}) ci-dessous &
  trivial (un log/exp par couple de technologies, jamais par
  transaction) \\[0.4em]
(c) exposants différents & $\gamma_\ell\neq\gamma_b$ &
  pas de forme fermée --- machinerie du rapport, \S3.2 &
  dépend du cache \\
\bottomrule
\end{tabular}
\end{center}

Formule fermée du régime (b) (rapport \S3, éq.14) :
\begin{equation}
\lambda^{*}=\frac{A_b^{1/(1-\gamma)}}{A_b^{1/(1-\gamma)}+A_\ell^{1/(1-\gamma)}}
\label{eq:lambda-star}
\end{equation}

\textbf{Point à ne pas manquer, cause probable de la confusion la plus
fréquente sur ce point} : le régime (b) n'est \textbf{pas} la formule
historique. $\lambda^{*}$ n'est égal à $1/2$ que si $A_\ell=A_b$ en plus
de $\gamma_\ell=\gamma_b$ (régime a) --- vérifier : à $A_b=A_\ell=A$,
$A_b^{1/(1-\gamma)}=A_\ell^{1/(1-\gamma)}$ et $\lambda^{*}=1/2$
exactement, ce qui redonne $\delta^{*}=(K_\ell-K_b)/2$ ; dès que
$A_b\neq A_\ell$, $\lambda^{*}\neq1/2$. Une intervention qui ne modifie
que $A$ (\S1, levier primaire) fait donc déjà sortir toute paire mixte du
régime (a) vers le régime (b), avec un partage \textbf{pondéré} du
capital --- « si les deux gammas sont égaux, le montant est aisément
calculable » (annotation de l'utilisateur) signifie « calculable en forme
fermée », pas « égal à la moitié de l'écart ».

\subsection*{3.2 Régime (c) : architecture hot/warm/cold du rapport}

Le rapport résout le régime (c) en exploitant une propriété structurelle
de M4.3Live, pas générique : $(A_i,\gamma_i)$ prend un petit nombre de
valeurs discrètes \textbf{exactes} (des « technologies »), pas des
perturbations continues (rapport \S6.1) --- la bonne opération est de
conditionner exactement sur le couple \textbf{ordonné} de technologies
(émettrice, receveuse), pas d'approximer dans un espace à 6 paramètres
continus (rapport \S6.3 : un surrogate générique à 6D construit avant
l'inspection du cas réel s'est avéré inutile une fois cette structure
exploitée --- ne pas répéter cette impasse). À couple de technologies
fixé, toute la non-linéarité tient dans une seule dimension $C=x+y$ ---
le Hessien de $\delta^{*}$ est exactement de rang un (rapport \S6,
éq.29-32), indépendamment des corrélations entre $x$ et $y$.

\textbf{Architecture exigée} (rapport \S11, « architecture recommandée »
--- suivre sauf justification écrite d'un meilleur choix dans le rapport
de conception, \S9) :

\begin{itemize}
\item \textbf{Indexation par technologie, pas par flottant.} Chaque
  $(A_i,\gamma_i)$ distinct reçoit un identifiant entier à sa première
  apparition ; une matrice dense
  \code{kernel[s\_émettrice][s\_receveuse]} route un couple en $O(1)$
  (rapport \S7.1).
\item \textbf{Chemin chaud} : régime (b) dès que
  $\gamma_{\text{émettrice}}=\gamma_{\text{receveuse}}$ --- aucune
  table, formule fermée (\S3.1).
\item \textbf{Chemin tiède} : couple de technologies nouveau, ou table
  en cours de construction --- noyau sans table (une étape de Newton
  depuis $z_0=-2t$, rapport \S4.2 éq.21, ou deux tours de point fixe) le
  temps d'observer assez d'occurrences.
\item \textbf{Chemin froid} : Newton robuste (rapport \S4, formulation
  exacte en $(t,u,z)$, éq.17) pour la première occurrence d'un couple, un
  $C$ hors du domaine couvert par la table, ou un audit aléatoire.
\item \textbf{Construction paresseuse d'une LUT 1D} $h(C)$ par couple de
  technologies non trivial, après un seuil d'amortissement \textbf{mesuré
  sur la machine cible} (indicatif du rapport : $\sim$1800 usages face à
  Newton exact sur son banc scalaire, \S8.1 --- \textbf{à recalibrer}, ce
  chiffre vient d'un benchmark synthétique, pas de ce moteur) ;
  interpolation linéaire ou cubique de Hermite selon la tolérance
  requise (rapport \S7.2 pour le pseudo-algorithme complet, à reprendre).
\item \textbf{Aucun tirage aléatoire pendant la construction/extension
  d'une table} --- la compilation peut s'exécuter en tout début du pas où
  une intervention est appliquée (même point de synchronisation que
  \S2) ; elle change la durée murale du pas, jamais la trajectoire
  (rapport \S7.2, dernier paragraphe).
\item \textbf{La LUT 2D globale n'est pas nécessaire} dans M4.3Live :
  elle resterait le baseline si les exposants redevenaient continus,
  mais le support discret transforme chaque cellule en problème 1D moins
  coûteux (rapport \S11, conclusion).
\end{itemize}

\textbf{Contrainte institutionnelle à fixer explicitement, pas seulement
numérique} (rapport \S2, éq.9) : le marché ne prête aujourd'hui que de
la plus riche vers la plus pauvre
(\code{K[lender] >= K[borrower]}, \code{m4\_3\_credit\_soc/m4\_3/model.py:
357-361}, comportement structurel hérité sans changement, \S3.3). Deux
choix possibles pour la borne supérieure de $\delta^{*}$ : (i) autoriser
le transfert jusqu'à l'optimum général $h(C)$, qui peut dépasser
l'égalisation des capitaux ; (ii) le plafonner à l'égalisation
($(K_\ell-K_b)/2$ dans un cas homogène, une borne équivalente en
général). Ce prompt ne tranche pas --- décision à documenter dans le
rapport de conception, et listée en \S11.

\subsection*{3.3 Hérité sans changement}

\begin{itemize}
\item Le sens du prêt (plus riche vers plus pauvre), le pool $k\equiv2$,
  les faillites cancel+destroy, l'ordre des phases --- tout hérité de
  M4.3 comme référence structurelle.
\item La formule de production $A\cdot K^{\gamma}$ par entité --- \S1.
\item $A_i$ et $\gamma_i$ par entité sont deux listes parallèles sur
  \code{Population}, fixées à la naissance, jamais retirées ni
  réutilisées après une mort --- même convention que les listes
  existantes de \code{Population} (\code{K}, \code{alive}, \code{birth},
  \ldots) dans le moteur de référence.
\item \code{\_pair\_rate} (\code{m4\_3\_credit\_soc/m4\_3/model.py:
  258-266}, $m=A\gamma K^{\gamma-1}$, $\mathit{rate}=\sqrt{m_\ell m_b}$)
  reste le mécanisme de taux \textbf{par défaut} de M4.3Live --- voir
  \S3.4, ce n'est pas remplacé par défaut, et sa généralisation par
  entité (chaque côté utilisant son propre $A_i,\gamma_i$) est directe
  et sans ambiguïté, contrairement au problème du principal.
\end{itemize}

\subsection*{3.4 Le taux d'intérêt : piste ouverte, pas un prérequis bloquant}

Une annotation de l'utilisateur propose d'aller plus loin que le
rapport : au lieu de garder \code{\_pair\_rate} inchangé, faire en sorte
que « les entités se partagent le surplus dû à leur coopération juste
après leur contrat », via un paramètre de partage $p$ qui « remplace
$m$ ». Le rapport, lui, ne couvre \textbf{que} le problème du montant
(son \S2 dit explicitement que remplacer le principal par ce transfert
optimal serait « une nouvelle règle de principal », distincte du calcul
du taux) --- cette extension au taux est une piste de l'utilisateur à
explorer, \textbf{pas un résultat déjà établi}, et ce prompt ne la
considère pas comme résolue.

\textbf{Ce qu'un candidat $p$ devrait produire, précisément.} Le carnet
de prêts a besoin d'un taux scalaire \code{rate} tel que
$\mathit{due} = \mathit{principal} \times \mathit{rate}$ à chaque pas,
perpétuellement (\code{m4\_3\_credit\_soc/m4\_3/model.py:166-187} pour
l'agrégat \code{due}, \code{model.py:572-594} pour le service payé chaque
pas --- mécanique à reproduire à l'identique dans le fork, \S3.5). Le
surplus coopératif, lui, est une quantité \textbf{ponctuelle}, calculée
une seule fois au moment du contrat :

$$\Delta = \bigl[A_b(x+\delta^{*})^{\gamma_b} +
A_\ell(y-\delta^{*})^{\gamma_\ell}\bigr] - \bigl[A_b x^{\gamma_b} +
A_\ell y^{\gamma_\ell}\bigr].$$

\textbf{Le problème réel n'est pas « comment partager $\Delta$ »} (un
$p\in[0,1]$ suffit à répondre à cette partie) --- \textbf{c'est que
transformer une quantité ponctuelle $\Delta$ en un taux perpétuel $r$
n'est pas un problème bien posé sans une hypothèse supplémentaire} (un
horizon implicite, un taux d'actualisation, une convention de rente
perpétuelle\ldots). C'est cette hypothèse manquante, pas le partage
lui-même, qui est la vraie question ouverte à trancher.

\textbf{Ce que ce prompt exige si cette piste est explorée} :

\begin{itemize}
\item Documenter explicitement l'hypothèse choisie pour transformer
  $(\Delta, p, \mathit{principal})$ en un taux $r$ --- pas un choix
  implicite enterré dans le code.
\item Vérifier que $r>0$ et que ce taux ne casse pas mécaniquement la
  solvabilité au pas suivant : \code{due} entre dans le calcul
  \textbf{avant} la dépréciation (\code{model.py:572-594}) --- un taux
  dérivé du surplus peut être très supérieur à l'ancien taux de
  rendement marginal, et un pic de faillites causé par cet artefact
  numérique n'aurait rien à voir avec la question du rebond (\S1).
  Documenter cette vérification, pas seulement l'effectuer.
\item \textbf{\code{\_pair\_rate} (\S3.3) reste le mécanisme
  opérationnel par défaut} tant que cette piste n'est pas terminée et
  validée par les deux points ci-dessus. La question du rebond (\S7) est
  mesurable avec le mécanisme de taux existant --- ne pas bloquer tout le
  programme sur $p$. Traiter ce point comme un chantier de conception
  \textbf{parallèle}, à documenter dans le rapport de conception (\S9)
  qu'il ait abouti ou non, pas comme un préalable au reste du prompt.
\end{itemize}

\subsection*{3.5 Fork indépendant, pas un import}

\textbf{M4.3Live doit être indépendant des autres modèles du dépôt}
(décision explicite de l'utilisateur) --- ce qui referme la question
laissée ouverte dans la révision précédente de ce prompt (« import
direct vs fork »). Ne \textbf{pas} importer \code{m4\_3\_credit\_soc.m4\_3}
comme dépendance : écrire le moteur de M4.3Live comme un paquet autonome
dans le nouveau dossier (\S10), qui lit
\code{m4\_3\_credit\_soc/m4\_3/model.py} comme référence structurelle
(mêmes conventions : \code{Population}/\code{LoanBook} à listes
parallèles jamais réindexées, RNG \code{numpy.random.default\_rng(seed)},
ordre des phases) sans en dépendre au sens Python du terme.

\subsection*{3.6 Parité avec M4.3 : souhaitable si simple, plus une exigence}

\textbf{Ce point change par rapport à la révision précédente de ce
prompt} (annotation explicite de l'utilisateur) : une correspondance
bit-à-bit avec M4.3 n'est \textbf{plus requise}. Si elle s'obtient
simplement dans le régime (a) (\S3.1) --- c'est-à-dire si le chemin
homogène du fork produit exactement $(K_\ell-K_b)/2$ par construction
plutôt que par coïncidence numérique --- c'est un plus à documenter, pas
un objectif à poursuivre activement. Les régimes (b) et (c) n'ont de
toute façon aucun analogue dans M4.3 littéral (qui n'implémentait que le
cas homogène) : aucune parité n'est même définie pour eux.

\section*{4. Architecture « en direct » : session, boucle de pas, journal d'interventions}
\addcontentsline{toc}{section}{4. Architecture « en direct » : session, boucle de pas, journal d'interventions}

\begin{itemize}
\item \textbf{Boucle de simulation} : un thread dédié appelle
  \code{Simulation.step()} en boucle à une cadence pilotable
  (lecture/pause, pas-à-pas, vitesse), avec une porte
  (\code{threading.Event} ou équivalent) pour la pause --- voir
  l'invariant RNG du \S2. La file d'attente d'interventions est un objet
  partagé protégé par verrou, purement stdlib
  (\code{threading.Lock}/\code{queue.Queue}), dans l'esprit de
  \code{simulation\_lab/jobs.py} (déjà \code{threading}-based) plutôt que
  le modèle \code{multiprocessing.Process} de
  \code{\_run\_single\_model\_worker} (\code{simulation\_lab/jobs.py:24-51})
  --- un sous-processus séparé rendrait le partage d'état (population
  potentiellement jusqu'à \code{pop\_max=30000} entités) coûteux à faire
  transiter par \code{Queue} à chaque pas ; documenter ce choix dans le
  rapport de conception.
\item \textbf{API de pilotage unique, partagée entre l'IHM HTML et un
  pilote sans tête} : la fonction qui applique une intervention (portée,
  paramètre, valeur, éventuel $t_0$/$\varphi$) doit être un point
  d'entrée Python appelable aussi bien depuis les gestionnaires HTTP
  \code{do\_POST} de l'IHM que depuis un script headless --- la campagne
  de mesure du \S7 nécessite plusieurs graines et un plan d'intervention
  identique, rejouable par script.
\item \textbf{Flux d'état vers le navigateur} : suivre la convention de
  sondage déjà en place dans \code{simulation\_lab}
  (\code{GET /api/jobs/<id>} interrogé périodiquement, cf.
  \code{simulation\_lab/web/static/app.js}) --- endpoint JSON renvoyant
  l'état courant, sondé à intervalle fixe côté client. Pas de nouvelle
  dépendance (pas de websocket tiers).
\item \textbf{Journal d'interventions, obligatoire, artefact du run} :
  chaque intervention effectivement appliquée ($t$ effectif, paramètre,
  portée, ancienne valeur, nouvelle valeur, et pour \code{fraction} les
  identifiants effectivement tirés) est écrite dans un fichier à côté
  des autres artefacts du run. \textbf{Propriété exigée et testée
  (\S8)} : rejouer un run en mode headless à partir de (seed, paramètres
  initiaux, journal d'interventions) doit reproduire une trajectoire
  strictement identique à l'originale.
\item \textbf{Statuts de fin de simulation}
  (\code{"extinction"}/\code{"explosion"}, \code{model.py:678-681}) : la
  boucle en direct doit s'arrêter proprement sur ces statuts et
  l'afficher clairement côté IHM.
\item \textbf{Reprise depuis un run M4.3 existant} (demande explicite de
  l'utilisateur) : pouvoir, à partir d'un run M4.3 déjà stocké dans
  \code{simulation\_lab\_data/runs/<run\_id>/} (paramètres et graine dans
  ses métadonnées, \S6), lancer une session M4.3Live à partir d'un pas
  $t_0$ arbitraire (1000, 2000, \ldots) de ce run. \textbf{Impossibilité
  à énoncer clairement, pas à contourner en silence} : M4.3Live étant un
  fork avec une institution différente (\S3), rejouer le run stocké avec
  le moteur M4.3Live ne reproduit sa trajectoire \textbf{que} dans le
  régime homogène (\S3.1a), et seulement si les chemins flottants
  concordent --- ce n'est pas garanti par construction (\S3.6). Exiger
  donc une \textbf{vérification de divergence obligatoire à chaque
  reprise} : rejouer depuis (seed, paramètres) jusqu'à $t_0$ avec le
  moteur M4.3Live, comparer la série obtenue à la série stockée du run
  original, rapporter l'écart maximal observé --- ne jamais présenter
  une reprise comme identique au run d'origine sans avoir mesuré cet
  écart.
\item \textbf{Snapshot complet à $t_0$, pour brancher plusieurs
  sessions} (demande explicite de l'utilisateur) : une fois un run amené
  à $t_0$ (reprise depuis M4.3 ci-dessus, ou session M4.3Live fraîche),
  sauvegarder un snapshot complet et picklable de l'état (population,
  carnet de prêts, RNG, caches du \S3.2) pour pouvoir relancer plusieurs
  sessions différentes (contrôle, un ou plusieurs traitements) depuis ce
  même $t_0$ sans re-simuler depuis $t=0$ à chaque fois --- sert
  directement le protocole apparié du \S7.
\end{itemize}

\section*{5. IHM HTML --- un outil de recherche interne, pas un produit}
\addcontentsline{toc}{section}{5. IHM HTML --- un outil de recherche interne, pas un produit}

Nouvelle page (ex. \code{/live}), servie par extension de
\code{simulation\_lab/web/app.py} ou par un serveur dédié isolé (à
trancher, \S11) mais dans les deux cas construite pour ne pas régresser
l'existant (\S6). Contenu minimal exigé :

\begin{itemize}
\item graphiques en direct des séries clés (\code{prod\_tot},
  \code{K\_tot}, \code{pop}, \code{defaults}, et une vue de la dispersion
  de $A$/$\gamma$ dans la population --- au minimum moyenne par cohorte
  née avant/après chaque intervention) ;
\item panneau de contrôle : paramètre, nouvelle valeur, portée
  (\code{all}/\code{new}/\code{fraction} avec champ $\varphi$ quand
  pertinent, \S2), avec les dégénérescences du \S2 explicitement
  affichées (pas de sélecteur qui ment sur ce qu'il fait) ;
\item contrôles de lecture : play/pause/pas-à-pas/vitesse ;
\item journal visible des interventions appliquées ($t$, paramètre,
  portée, valeurs).
\end{itemize}

C'est un outil de recherche interne, que l'utilisateur veut pouvoir
manipuler lui-même en complément d'un travail autonome de l'agent
(remarque explicite de l'utilisateur) : priorité à la clarté et à
l'exactitude des chiffres affichés sur le raffinement visuel. Réutiliser
le style existant (\code{simulation\_lab/web/static/app.css}) plutôt
qu'en construire un nouveau. Ne pas sur-investir en design.

\section*{6. Intégration avec simulation\_lab : isolation, non-régression}
\addcontentsline{toc}{section}{6. Intégration avec simulation\_lab : isolation, non-régression}

\code{simulation\_lab/} reste fonctionnel pour ses usages actuels
(lancement simple/batch, \code{/launch}, \code{/results}, CLI
\code{list-models}/\code{run}/\code{batch}/\code{list-runs}) après ce
travail --- vérifier au sens propre (relancer ces commandes, pas
seulement lire le code) avant de considérer la tâche terminée. Isoler le
nouveau code plutôt que modifier en place les chemins existants de
\code{contracts.py}/\code{jobs.py}/\code{runs/} --- par exemple un
module \code{simulation\_lab/live/} séparé, avec ses propres routes
(\code{/live}, \code{/api/live/*}) ajoutées à
\code{SimulationLabHandler} (\code{simulation\_lab/web/app.py:45-193})
sans toucher aux branches existantes de \code{do\_GET}/\code{do\_POST}.
La fonctionnalité de reprise (\S4) lit les runs déjà stockés via
\code{simulation\_lab.runs.storage.RunStorage} (métadonnées, seed,
paramètres) --- en lecture seule, sans modifier le format de stockage
existant. Si un adaptateur \code{BaseSimulationModel} (\code{contracts.py})
a un sens pour exposer une version \emph{batch} (sans direct) du moteur
étendu via \code{/launch} classique, c'est un bonus optionnel --- pas un
prérequis, le pilote headless du \S4 suffit pour la science.

\section*{7. Protocole expérimental pour la question du rebond}
\addcontentsline{toc}{section}{7. Protocole expérimental pour la question du rebond}

\begin{itemize}
\item \textbf{Contrôle apparié par graine.} Deux runs de même graine,
  mêmes paramètres initiaux, l'un avec un journal d'interventions vide
  (contrôle), l'autre avec une intervention à $t_0$ (traitement) : les
  deux trajectoires sont \textbf{rigoureusement identiques jusqu'à
  $t_0$} (même séquence RNG consommée dans le même ordre jusque-là) et
  ne divergent qu'à partir de l'intervention. C'est ce qui rend la
  différenciation appariée légitime et peu coûteuse --- s'appuyer dessus
  comme estimateur plutôt qu'en inventer un autre. \textbf{Utiliser le
  snapshot à $t_0$ du \S4} pour brancher contrôle et traitement(s) sans
  re-simuler depuis $t=0$ à chaque fois.
\item \textbf{Plan minimal} : à $t_0$ choisi après relaxation (temps
  caractéristique du système, pas une valeur arbitraire). \textbf{Ne pas
  relancer de mesure de temps de relaxation à partir de zéro} : M4.3 a
  déjà estimé ce temps à sa baseline
  (\code{m4\_3\_credit\_soc/scripts/renewal\_relaxation\_all\_runs.py},
  résultats cités dans son rapport final) --- partir de cette estimation
  existante pour fixer $t_0$ (avec une marge de sécurité raisonnable), et
  ne remesurer que si les paramètres retenus pour \S7 s'écartent trop de
  la baseline M4.3 pour que l'estimation reste fiable. Comparer ensuite
  (a) contrôle, (b) portée \code{fraction} (ou \code{new}, \S1) à
  $\varphi$ fixé sur $A$ (levier primaire, \S1, régime (b) du \S3.1),
  (c) portée \code{all} de même amplitude sur $A$ comme référence
  « hausse globale », (d) au moins une variante faisant varier $\gamma$
  pour placer la paire dans le régime (c) du \S3.1 --- ceci exerce
  délibérément le chemin coûteux de l'institution, pas seulement le
  chemin fermé. Plusieurs graines par cellule (le nombre exact à ta
  discrétion, mais jamais une seule). Optionnellement, des jeux
  différents de $\lambda$/$\delta$/$\sigma$ peuvent servir de régimes de
  population/marché secondaires (\S2).
\item \textbf{Fenêtre d'observation} : mesurer \code{prod\_tot} sur une
  fenêtre post-$t_0$ dont la longueur est justifiée par le temps de
  relaxation du système, pas fixée arbitrairement.
\item \textbf{Verdict à rendre}, avec la discipline
  fait/inférence/hypothèse/incertitude du reste du dépôt : effet rebond
  observé / non observé / non tranchable avec le budget de calcul
  disponible (\S1 --- l'objectif de cette itération est l'observation,
  pas l'attribution causale).
\item Réutiliser l'outillage statistique existant du dépôt
  (\code{m4\_3\_credit\_soc/scripts/lib\_metrics.py} et voisins) si
  pertinent plutôt que le récrire ; le dire si rien ne s'applique.
\end{itemize}

\section*{8. Tests exigés}
\addcontentsline{toc}{section}{8. Tests exigés}

\begin{itemize}
\item \textbf{Institution de principal} (\S3) : régime (a) reproduit
  $(K_\ell-K_b)/2$ à une tolérance documentée (pas nécessairement
  bit-exacte, \S3.6) ; régime (b) vérifié exactement contre la formule
  fermée (rapport éq.14) ; régime (c) vérifié par comparaison
  Newton-exact vs le chemin choisi (chaud/tiède/froid), erreur sous la
  tolérance documentée.
\item \textbf{Rejeu/déterminisme} (\S4) : le journal source de ce test
  doit provenir d'une \textbf{session en direct réelle}, où les
  interventions (au moins une par portée) ont été soumises de façon
  asynchrone pendant que la boucle tournait --- pas d'un plan headless
  écrit et exécuté par le même code que celui qui doit être vérifié. Ce
  journal, rejoué en mode headless, doit reproduire une trajectoire
  strictement identique à l'originale, en particulier même $t$ effectif
  que celui enregistré.
\item \textbf{Divergence à la reprise} (\S4) : pour la fonctionnalité de
  reprise depuis un run M4.3 stocké, comparer systématiquement la série
  rejouée à la série originale et vérifier que l'écart mesuré est
  rapporté, pas seulement calculé en silence.
\item \textbf{Invariant de pause} (\S2) : $N$ pas avec une ou plusieurs
  pauses au milieu produisent la même trajectoire que $N$ pas sans
  pause.
\item \textbf{Sémantique de portée}, au moins un test unitaire par
  portée : \code{new} ne change rien aux entités déjà nées ; \code{all}
  change bien toutes les entités vivantes au pas suivant ;
  \code{fraction} ne touche que le nombre attendu d'entités \textbf{et}
  ne change pas le défaut utilisé par les naissances postérieures à
  $t_0$ (vérifier explicitement qu'une entité née après $t_0$ reçoit
  l'ancien défaut, pas la valeur tirée).
\item \textbf{Solvabilité}, si la piste du \S3.4 est implémentée : le
  taux candidat $r(p,\Delta,\mathit{principal})$ reste positif et ne
  provoque pas un pic de défauts mécanique au pas suivant, sur un
  échantillon de paires représentatif.
\item Assertions Python simples, cohérent avec la convention du dépôt
  (pas de pytest, \code{CLAUDE.md}).
\end{itemize}

\section*{9. Livrables}
\addcontentsline{toc}{section}{9. Livrables}

\begin{itemize}
\item Code complet et fonctionnel dans \code{m4\_3live\_credit\_soc/}
  (\S10 pour l'arborescence attendue) --- un système qui tourne
  réellement, pas un squelette : démarrer le serveur, piloter une
  session en direct dans un vrai navigateur, vérifier visuellement que
  les graphiques et le panneau de contrôle fonctionnent, avant de
  déclarer la tâche terminée.
\item \textbf{Rapport de conception}, \LaTeX{} compilé en PDF
  (\code{report/conception\_m4\_3live.pdf}) : toutes les décisions
  d'architecture de ce prompt justifiées avec citations
  \code{fichier.py:ligne} --- en particulier les trois régimes du \S3.1,
  l'architecture hot/warm/cold retenue (\S3.2) et ses seuils
  recalibrés, le choix de borne institutionnelle (\S3.2, dernier
  point), le statut du chantier taux/$p$ (\S3.4, abouti ou non), le
  caractère indépendant du fork (\S3.5), et les résultats de la
  vérification de divergence à la reprise (\S4) --- y compris les
  décisions explicitement laissées ouvertes (\S11) et la façon dont tu
  les as tranchées.
\item \textbf{Rapport de résultats}, \LaTeX{} compilé en PDF
  (\code{report/rapport\_final.pdf}) : protocole du \S7, résultats,
  figures, verdict d'observation sur l'effet rebond, limites --- même
  discipline scientifique que
  \code{m4\_3\_credit\_soc/report/rapport\_final.tex}.
\item Compilation PDF vérifiée (le \code{.tex} seul ne suffit pas).
\item Un \code{README.md} et un \code{JOURNAL.md} courts dans
  \code{m4\_3live\_credit\_soc/}, cohérents avec la convention du dépôt.
\end{itemize}

\section*{10. Arborescence attendue}
\addcontentsline{toc}{section}{10. Arborescence attendue}

\begin{verbatim}
m4_3live_credit_soc/
+-- prompts/PROMPT_M4_3LIVE.md      (ce document, source markdown)
+-- prompts/PROMPT_M4_3LIVE.tex/.pdf (cette version LaTeX/PDF)
+-- prompts/rapport_architecture_offline_online.pdf (etude du S3, lecture obligatoire)
+-- m4_3live/                        (paquet moteur INDEPENDANT, S3.5)
|   +-- model.py                     (institution du S3, per-entite A/gamma)
|   +-- live.py                      (session, boucle, journal, S4)
+-- driver/                          (pilote headless, S4 -- nom a ta discretion)
+-- web/ ou extension de simulation_lab/live/ (S6, a trancher S11)
+-- tests/
|   +-- test_institution.py          (S8, trois regimes)
|   +-- test_replay.py
|   +-- test_scopes.py
|   +-- test_resume_divergence.py    (S4, S8)
+-- scripts/                         (campagne du S7)
+-- results/                         (runs de la campagne, snapshots)
+-- report/
|   +-- conception_m4_3live.tex/.pdf
|   +-- rapport_final.tex/.pdf
|   +-- figures/
+-- README.md
+-- JOURNAL.md
\end{verbatim}

\section*{11. Décisions laissées ouvertes, à trancher et documenter toi-même}
\addcontentsline{toc}{section}{11. Décisions laissées ouvertes, à trancher et documenter toi-même}

\begin{itemize}
\item Nom exact du paquet/dossier si \code{m4\_3live\_credit\_soc}/
  \code{m4\_3live} s'avère gênant en pratique --- réglé par défaut,
  changement possible si justifié.
\item Extension de \code{simulation\_lab/web/app.py} en place (isolée,
  \S6) vs serveur HTTP dédié entièrement séparé pour \code{/live} --- les
  deux sont acceptables, documenter le choix et pourquoi.
\item Mécanisme exact de sondage HTTP (intervalle, format JSON) ---
  reprendre la convention \code{app.js} existante par défaut, l'adapter
  si le débit d'un direct l'exige, en le disant.
\item Étendre l'hétérogénéité par entité à d'autres paramètres que
  $A$/$\gamma$ (ex. \code{delta} par entité) --- hors mandat par défaut,
  seulement si le protocole du \S7 en révèle un besoin net.
\item \textbf{Co-variation de \code{K0} avec une intervention
  $A$/$\gamma$} (\S2) --- automatique, optionnelle sur demande, ou
  délibérément découplée ; à trancher avant la campagne du \S7 et à
  documenter, quelle que soit la réponse.
\item \textbf{Borne supérieure du transfert optimal} (\S3.2, égalisation
  des capitaux vs optimum général $h(C)$) --- choix institutionnel, pas
  numérique.
\item \textbf{Seuils hot/warm/cold} (\S3.2) --- recalibrer sur la
  machine cible, ne pas reprendre tels quels les chiffres du banc
  synthétique du rapport.
\item \textbf{Statut du chantier taux/$p$} (\S3.4) --- abouti (avec
  l'hypothèse de conversion $\Delta\to r$ documentée et validée) ou
  explicitement ajourné au profit de \code{\_pair\_rate}, mais pas
  laissé dans un état intermédiaire non documenté.
\end{itemize}

\section*{12. Consignes de méthode de travail, pour toi (Opus 5)}
\addcontentsline{toc}{section}{12. Consignes de méthode de travail, pour toi (Opus 5)}

\begin{itemize}
\item \textbf{Livrer un système complet, pas des ébauches.} Chaque pièce
  du \S10 doit fonctionner de bout en bout avant que la tâche soit
  déclarée terminée --- pas de squelette avec des TODO, pas de
  fonctionnalité à moitié câblée.
\item \textbf{Calibrer la longueur des livrables écrits sur leur
  substance.} Les deux rapports \LaTeX{} doivent couvrir ce que ce
  prompt demande, sans sections de remplissage, sans résumé redondant en
  fin de document, sans paragraphe qui répète ce qui vient d'être dit
  autrement.
\item \textbf{Pas de vérification en boucle au-delà de ce qui est
  demandé.} Les tests du \S8 et le test manuel dans un vrai navigateur
  (\S9) sont les vérifications exigées --- inutile d'ajouter des passes
  de relecture supplémentaires au-delà.
\item \textbf{Périmètre : livrer ce qui est demandé, à l'échelle
  demandée.} Prendre seul les décisions d'implémentation de routine (y
  compris celles listées en \S11) ; ne pas élargir le mandat
  scientifique au-delà du \S7 ; si une lecture différente de ce prompt
  changerait matériellement le travail à produire, le signaler en une
  phrase et continuer avec la lecture la plus proche du texte plutôt que
  de redéfinir la tâche.
\end{itemize}

\subsection*{Sous-agents : subdiviser intelligemment, pas systématiquement}
\addcontentsline{toc}{subsection}{Sous-agents : subdiviser intelligemment, pas systématiquement}

Ce dépôt et cette tâche se prêtent à de la délégation --- utiliser les
sous-agents avec discernement plutôt que tout faire en séquence toi-même
ou, à l'inverse, tout déléguer par réflexe :

\begin{itemize}
\item \textbf{Réserver ton propre raisonnement (Opus) aux décisions qui
  déterminent la validité scientifique et architecturale du travail} :
  l'institution de principal et ses trois régimes (\S3), la sémantique
  exacte des trois portées (\S2), le protocole expérimental et
  l'écriture du verdict (\S7), l'intégration avec \code{simulation\_lab}
  sans régression (\S6), et la relecture finale de cohérence entre code
  et rapports. Ce sont les endroits où une erreur de jugement invalide
  tout le reste.
\item \textbf{Déléguer les missions simples, bien délimitées et
  vérifiables isolément à des sous-agents sur un modèle plus léger
  (Sonnet ou Haiku)} : écrire le CSS/JS répétitif de l'IHM une fois la
  structure de l'API fixée, écrire un test unitaire donné une
  spécification précise (\S8), lancer et collecter une graine de la
  campagne du \S7 une fois le protocole gelé, rédiger une section
  descriptive du rapport de conception à partir de décisions déjà
  prises et citées, exécuter une vérification de non-régression sur
  \code{simulation\_lab} et rapporter le résultat brut.
\item Donner à chaque sous-agent délégué un mandat autonome et complet
  (pas un fragment de phrase) : quel fichier, quel contrat
  d'entrée/sortie, quel test de validation --- puisqu'un sous-agent frais
  ne voit pas ce prompt, lui fournir directement les extraits pertinents
  (numéros de ligne, formules, sémantique de portée) plutôt que d'y
  renvoyer par référence.
\item Ne pas déléguer la vérification de ton propre travail à un autre
  sous-agent en plus des tests déjà exigés (\S8/\S9) --- pas de double
  passe de relecture par sous-agent interposé.
\item Ne pas ouvrir plus de sous-agents que la structure du travail n'en
  justifie --- une tâche qui tient en quelques appels d'outils directs
  n'a pas besoin d'être déléguée.
\end{itemize}

\end{document}
